\begin{abstract}
Protecting electronic health record (EHR) telemetry and sensitive patient data in distributed healthcare systems requires computationally efficient, reversible block encryption schemes. This paper presents a \textbf{SHA-512-Based 256-bit Block Cellular Automata Cryptosystem} tailored for secure healthcare applications. The proposed architecture derives a 256-bit binary key ($K_{256}$) from the first 32 bytes of a SHA-512 password digest. Input text is framed with a 32-bit binary length header, zero-padded, and processed in 256-bit blocks mapped to $16 \times 16$ binary matrices. Non-linear diffusion is achieved using a reversible $2 \times 2$ Margolus Cellular Automata (CA) transformation featuring a 4-bit local bijective rule table, alternating standard (odd rounds) and shifted (even rounds) periodic partitions, and a 2D circular matrix shift (down by 2 rows, right by 2 columns) after each round. Evaluated across round counts $R \in \{2, 4, 6, 8, 10\}$, empirical metrics demonstrate that diffusion scales monotonically with additional CA rounds. For the proposed 10-round configuration, the transformation achieves a mean bit-change count of \textbf{81.92 / 256} bits (\textbf{32.00\%} bit avalanche), a Byte NPCR of \textbf{84.62\%}, a Byte UACI of \textbf{22.83\%}, and an average execution latency of \textbf{1.883 ms} per 256-bit block. Across 100 deterministic synthetic healthcare records, the cryptosystem demonstrates \textbf{100\%} loss-less reversibility and \textbf{100\%} exact plaintext record recovery. The 10-round configuration provides a practical balance between enhanced non-linear bit diffusion and computational efficiency for healthcare data protection.
\end{abstract}

\begin{IEEEkeywords}
Cellular Automata, Block Cipher, Margolus Neighborhood, Key Derivation, SHA-512, Strict Avalanche Criterion, Healthcare Security, Performance Benchmarking.
\end{IEEEkeywords}
